\documentclass[a4paper, 11pt]{article}
\usepackage[T1]{fontenc} 
\usepackage[english]{babel}
\usepackage[utf8]{inputenc}
\usepackage[top=3cm, bottom=3cm, left = 2cm, right = 2cm]{geometry} 
\geometry{a4paper} 
\usepackage[utf8]{inputenc}
\usepackage{textcomp}
\usepackage{graphicx} 
\usepackage{amsmath,amssymb}
\usepackage{booktabs}
\usepackage{bm}  
\usepackage[pdftex,bookmarks,colorlinks,breaklinks]{hyperref}  
\hypersetup{linkcolor=blue,citecolor=blue,filecolor=blue,urlcolor=blue} 
\usepackage{memhfixc} 
\usepackage{pdfsync}  
\usepackage{fancyhdr}
\pagestyle{fancy}

\title{Sample Report}
\author{Author McWriterson}

\begin{document}
\maketitle
\begin{center}\small English translation of the Polish original. Historical results and formulas retained.\end{center}
\tableofcontents

\section{Task 1}
\subsection{Task:}
\noindent \textbf{Given:} \( x \), \( e \): column vectors in \( \mathbb{R}^n \).

\[
k = \frac{\|x\|}{\|e\|} = \sqrt{\frac{x^T x}{e^T e}}, \quad v = x - k e, \quad u = x + k e.
\]
Write a function
\begin{center}
\begin{verbatim}
function [P, Q] = haus(x, e)
\end{verbatim}
\end{center}
returning matrices \( P \) and \( Q \) given by:
\[
P = I - 2 \frac{v v^T}{v^T v}, \quad Q = I - 2 \frac{u u^T}{u^T u}.
\]

Function call:
\begin{center}
\begin{verbatim}
x = [1; 0; 0; 0];
e = [1/2; -1/2; -1/2; -1/2];
[P, Q] = haus(x, e)
\end{verbatim}
\end{center}
< from the slides >
\subsection{Questions:} 

\subsection{Results:} 
\begin{center}
\begin{tabular}{@{}|l|l|l|l|l|l|@{}}
\toprule
n                         & k      & Error 1      & Error 2      & Error 3      & Error 4      \\ \hline
3                         &$ 1.0892$ & $2.2204 e-16$ & $2.2204 e-16$ & $2.2204 e-16$ & $2.2204 e-16$ \\ \hline
3000                      & $0.9829$ & $1.7764 e-15$ & $1.9984 e-15$ & $3.2198 e-15$ & $2.9976 e-15$ \\ \hline
3000; e almost equal to x &$ 1.000$ & $1.9984 e-15$ & $1.8874 e-15$ &$ 2.6645 e-15$ & $3.8858 e-15$ \\ \hline
3000; x, e orthogonal: & $0.9939$ & $1.9984 e-15$ & $1.7764 e-15$ & $4.2188 e-15$ & $3.5527 e-15$ \\ \hline
\end{tabular}
\end{center}
\subsection{Conclusions:}

\section{Task 2}
\subsection{Task:}
\noindent \textbf{Given:} \( B := A \), a nonsingular matrix.

\begin{center} \textbf{QR algorithm:} \end{center}
\begin{algindent}
    \item \textbf{for} \( m := 1 \) \textbf{to} \( n \) \textbf{do}
    \item \hspace{1.5em} \( t := \sum_{s=m}^n (B[s, m])^2 \);
    \item \hspace{1.5em} \textbf{if} \( B[m, m] < 0 \) \textbf{then} \( k := \sqrt{t} \) \textbf{else} \( k := -\sqrt{t} \);
    \item \hspace{1.5em} \( d[m] := k \);
    \item \hspace{1.5em} \( f[m] := \frac{1}{k \cdot B[m, m] - t} \);
    \item \hspace{1.5em} \( B[m, m] := B[m, m] - k \);
    \item \hspace{1.5em} \textbf{for} \( i := m+1 \) \textbf{to} \( n \) \textbf{do}
    \item \hspace{3em} \( y := f[m] \sum_{s=m}^n B[s, m] B[s, i] \);
    \item \hspace{3em} \textbf{for} \( j := m \) \textbf{to} \( n \) \textbf{do} \( B[j, i] := B[j, i] + B[j, m] \cdot y \);
\end{algindent}

\noindent Result:
\[
R = \begin{bmatrix}
d[1] & B[1, 2] & \cdots & B[1, n] \\
0 & d[2] & B[2, 3] & \cdots & B[2, n] \\
\vdots & \ddots & \ddots & \cdots \\
0 & \cdots & d[n-1] & B[n-1, n] \\
0 & \cdots & 0 & d[n]
\end{bmatrix}, \quad Q = A R^{-1}.
\]

Write a function:
\begin{center}
\begin{verbatim}
function [Q, R] = QR(A)
\end{verbatim}
\end{center}
returning \( Q \) and \( R \), then find QR decompositions of randomly sampled matrices.

\subsection{Questions:} 

\subsection{Results:} 
\begin{center}
\begin{tabular}{|l|l|l|l|l|}
\hline
n   & error-1       & error-2       & error-3       & error-4       \\ \hline
3   & $8.8818e-16$ & $5.5511e-17$ & $4.4409e-16$ & $3.3307e-16$ \\ \hline
500 & $5.4368e-13$ & $3.5416e-14$ & $1.1102e-15$ & $7.9381e-15$ \\ \hline
100 & $3.9968e-15$ & $5.6843e-14$ & $1.7764e-15$ & $3.1264e-13$ \\ \hline
\end{tabular}
\end{center}
\subsection{Conclusions:}

\section{Task 3}
\subsection{Task:}
For the matrix
\[
A = \begin{bmatrix}
a_{1,1} & a_{1,2} & \cdots & a_{1,n} \\
a_{2,1} & a_{2,2} & \cdots & a_{2,n} \\
\vdots & \vdots & \ddots & \vdots \\
a_{n,1} & a_{n,2} & \cdots & a_{n,n}
\end{bmatrix}
\]

The matrix 2-norm of \( A \) is the maximum over the roots of the characteristic polynomial of \( A^T A \), that is:
\[
w(x) := \det(A^T A - x I), \quad C := \{ x : w(x) = 0 \}, \quad \|A\|_2 = \max_{x \in C} \sqrt{x}.
\]

Write a function
\begin{center}
\begin{verbatim}
ret = norm2(A)
\end{verbatim}
\end{center}
returning the matrix 2-norm of \( A \).
\subsection{Questions:} 

\subsection{Results:} 

\subsection{Conclusions:}
Comparing our $norm2$ function with Octave's built-in $norm$, we obtain discrepancies due to numerical errors. Our function's errors mainly come from computing characteristic-polynomial roots, an operation particularly sensitive to closely spaced eigenvalues. Complex roots should not occur; however, for flat matrices the errors may be complex because the polynomial roots have small imaginary parts. Rounding during characteristic-polynomial computation can also produce unwanted complex roots for large matrices. Large diagonal values can lead to incorrect rounding caused by memory overflow at extreme values. To avoid root-computation errors, one can use the power method. It focuses on the largest eigenvalue, which becomes more apparent across iterations, and requires only matrix--vector multiplication, reducing the risk of errors.




\section{Task 4}
\subsection{Iterative methods}
\subsection{Task}

\subsubsection{Jacobi method}
\noindent \textbf{Given:} \( A \): an \( n \times n \) matrix; \( b \), \( x_0 \): column vectors in \( \mathbb{R}^n \); \( k \): a natural number.

Let \( B := D \).

Write a function
\begin{center}
\begin{verbatim}
function retval = Jacobi(A, b, x0, k)
\end{verbatim}
\end{center}
returning the vector \( x^{(k)} \).

\subsubsection{Gauss--Seidel method:}

\noindent \textbf{Given:} \( A \): an \( n \times n \) matrix; \( b \), \( x_0 \): column vectors in \( \mathbb{R}^n \); \( k \): a natural number.

Let \( B := D + L \).

Write a function
\begin{center}
\begin{verbatim}
function retval = GaussSiedel(A, b, x0, k)
\end{verbatim}
\end{center}
returning the vector \( x^{(k)} \).

\subsubsection{SOR method:}

\noindent \textbf{Given:} \( A \): an \( n \times n \) matrix; \( b \), \( x_0 \): column vectors in \( \mathbb{R}^n \); \( k \): a natural number; \( t \in (0, 2) \).

Let \( B := \frac{1}{t} D + L \).

Write a function
\begin{center}
\begin{verbatim}
function retval = SOR(A, b, x0, k, t)
\end{verbatim}
\end{center}
returning the vector \( x^{(k)} \).


\subsection{Results and conclusions:} 
For all calculations we set $x=0, b=rand(n,1) *10$, perform $1000$ trials and average their results.\\
\subsubsection{Parameter $k$}
We estimate \( k \), the number of iterations needed to establish convergence. To avoid errors due to a lack of diagonal dominance, set \( c = n \cdot 10 \). This gives a dominant diagonal and removes errors arising from matrix structure. We can then set \( n = 10 \), because this matrix size avoids size-related and convergence problems.

\begin{center}
\begin{tabular}{|l|l|l|l|l|l|}
\hline
Method       & 1          & 3          & 5          & 10         & 16         \\ \hline
Jacobi       & 4.0556e-03 & 6.5079e-06 & 9.9949e-09 & 2.7341e-15 & 1.1819e-17 \\ \hline
Gauss-Siedel & 3.3702e-03 & 4.4534e-07 & 4.7716e-12 & 2.6600e-17 & 1.8228e-17 \\ \hline
SOR          & 0.010987   & 1.3029e-04 & 1.2422e-06 & 2.5774e-11 & 2.3522e-17 \\ \hline
\end{tabular}
\end{center}
Conclusions:
\begin{itemize}
    \item SOR: effect of a poorly chosen relaxation factor:\\
    Although SOR should theoretically converge fastest, our calculations do not confirm this. A fixed, poorly chosen relaxation factor was used, reducing its efficiency. A poor relaxation factor can therefore substantially reduce the method's effectiveness.
    \item Comparison of Jacobi and Gauss--Seidel:\\
    The results show that Gauss--Seidel converges faster from the first iteration because new values are immediately reused within it. Jacobi uses only preceding-iteration values, slowing convergence.
\end{itemize}



\subsubsection{Parameter $c$}
We estimate \( c \), the value added to the diagonal in \( A + c \cdot \text{eye}(n) \). Set \( k = 100 \) to ensure convergence.
\begin{itemize}
    \item Case $c = n/2$
\begin{center}
\begin{tabular}{|l|l|l|l|}
\hline
Method       & $3$            & $10$           & $100$          \\ \hline
Jacobi       & $1.4872e-08$ & $2.0980e-03$ & $8.8618e-03$ \\ \hline
Gauss-Siedel & $5.8098e-16$ & $3.9518e-16$ & $17138e-16$  \\ \hline
SOR          & $5.9233e-16$ & $3.9828e-16$ & $17124e-16$  \\ \hline
\end{tabular}
\end{center}
    \item Case $c = 10*n$
    \begin{center}
\begin{tabular}{|l|l|l|l|}
\hline
Method       & 3            & 10           & 100          \\ \hline
Jacobi       & $3.9344e-17$ & $2.5617e-17$ & $1.0956e-17$ \\ \hline
Gauss-Siedel & $4.0613e-17$ & $2.5840e-17$ & $1.0976e-17$ \\ \hline
SOR          & $3.7866e-17$ & $2.5573e-17$ & $1.0902e-17$ \\ \hline
\end{tabular}
    \end{center}
\end{itemize}
Conclusions:\\
The approximation of $c$ suggested by the mean $E[\sum_i^n X_i]$, where $X_1 \sim \mathrm{U}([0, 1])$, appears to be a good estimate. The Jacobi errors must arise from the algorithm itself, but I do not understand why they occur for [sentence unfinished in the original].
\subsection{Suggestions}
Stopping criterion:\\
Choose $p$ and define a stopping time at step $n$ when the following condition holds:
$$\left|\frac{x_n - x_{n-1}}{x_n}\right| \times 100 < p$$




\end{document}
